\section{Uniform certificate tools}
\label{p14:cert:section}

Fixed-order estimates leave a separate problem: summing orders that grow
with $N$. The arithmetic certificates below solve it by counting
congruences and graph components with explicit cardinality weights.
Their definitions concern the actual values in the windows and apply
before any affine prescription is imposed. In particular, the geometric
counts also permit irregular windows; replacing $p$ by $p_x\le2p$ only
changes a fixed weight. The final parameter choices for the optimal
reconstruction are made in Section~\ref{p13:transfer:section}.
The complete elementary proofs are collected in
Section~\ref{p14:comp:certificate-tools}.

\subsection{Small kernels and moderate congruence certificates}

For $Y\ge B$, write
\[
 K_Y(v)=\prod_{\substack{\ell>Y\ \mathrm{prime}\\v_\ell(v)\ \mathrm{odd}}}\ell.
\]
The following elementary form of the Rankin argument will be used with
its exponent fixed.

\begin{lemma}[Counting a small truncated squarefree kernel]
\label{p1:core:rankin}
For fixed $1/2<\sigma<1$, fixed $C_1>0$, and $T\ge1$,
\begin{equation}
 \#\{v\le C_1N:K_Y(v)\le T\}
 \ll_{\sigma,C_1}N^\sigma T^{1-\sigma}
       \exp(C_\sigma Y^{1-\sigma}).
 \label{p1:core:rankin-bound}
\end{equation}
\end{lemma}
\noindent The complete proof is given in Section~\ref{p14:comp:certificate-rankin}.


Fix
\begin{equation}
 A_0>1,\quad 0<\beta<1,\quad \tfrac12<\sigma<1,
 \quad A_0(1-\sigma)<1,
 \label{p1:core:parameter-conditions}
\end{equation}
and put $H=\log N$, $Y=H^{A_0}$, $T=N^\beta$ and
$\delta=(1-\sigma)(1-\beta)>0$. Mark a window if one of its values
has $K_Y(v)\le T$. Each value lies in at most $B$ windows, so the
number $M$ of marked centres satisfies
\begin{equation}
 M\le N^{1-\delta+o(1)}.
 \label{p1:core:marks}
\end{equation}
Here and below the $o(1)$ depends on the fixed parameters, but not on
the size of a family subsequently summed over.

Choose two distinct values in an unmarked anchor window. Their kernels
are coprime, because a prime exceeding $Y\ge B$ cannot divide their
nonzero difference. Thus their product $Q$ is squarefree and
$T^2<Q\le9N^2$. If a relation selects these values, each prime of $Q$
occurs to odd order in at least one other window. Assign it to one such
occurrence. For each receiving block $j$, let $D_j$ be the product of
the assigned primes. The $D_j$ are coprime and $\prod_jD_j=Q$.
Call the resulting certificate \emph{moderate} when every $D_j\le6NB$.
After specifying a certificate we may discard the relation itself:
the congruences alone bound the possible centres.

\begin{lemma}[Counting moderate certificates]
\label{p1:core:moderate}
Let $\mathcal U_r$ be the ordered lists of $r$ original windows having
such a certificate, and set
$m=\lceil 2\log(3N)/\log Y\rceil$. Then
\begin{equation}
 |\mathcal U_r|
 \le36rB^4 2^{r-1}N^rT^{-2}(rB)^m.
 \label{p1:core:moderate-count}
\end{equation}
For every fixed $u>0$,
\begin{equation}
 \sum_{r\ge1}\frac{(up)^r}{r!}|\mathcal U_r|
 \le N^{-2\beta+4/A_0+o(1)}.
 \label{p1:core:moderate-baseline}
\end{equation}
\end{lemma}
\noindent The complete proof is given in Section~\ref{p14:comp:certificate-moderate}.


\subsection{Saturation has small degree in both directions}

Put a directed edge $x\to y$ when two values in window $x$ have a
kernel product possessing a divisor $D>6NB$ all of whose primes divide
values in window $y$. We impose no direct-sum condition on $S_x,S_y$.
Let $\Gamma$ be the underlying simple undirected graph.

\begin{lemma}[Incoming and outgoing saturation degrees]
\label{p1:core:degrees}
Uniformly in the fixed vertex,
\begin{equation}
 \deg^+(x),\deg^-(x)\le N^{2/A_0+o(1)},
 \qquad \Delta(\Gamma)\le N^{2/A_0+o(1)}.
 \label{p1:core:degree-bound}
\end{equation}
\end{lemma}
\noindent The complete proof is given in Section~\ref{p14:comp:certificate-degrees}.


If a primitive family has no moderate certificate, every unmarked
vertex has a neighbour in its induced saturation graph. To see this,
choose a nonzero relation, select two of its values at that vertex and
assign the primes to exterior occurrences. Some receiver must have
$D_j>6NB$. The choices made at different vertices need not agree.
Their existence in one deterministic graph is sufficient.

\begin{lemma}[A generating function with marked isolated vertices]
\label{p1:core:generator}
Let a simple graph have $n$ vertices, maximum degree $\Delta$ and
$M$ marked vertices. If $\mathcal F$ consists of the vertex sets
whose induced isolated vertices are all marked, then, for $z\ge0$,
\begin{equation}
 \sum_{D\in\mathcal F}z^{|D|}
 \le\exp\left\{Mz+nz\bigl((1+z)^\Delta-1\bigr)\right\}.
 \label{p1:core:generator-bound}
\end{equation}
More generally, for nonnegative vertex weights $w_v$ and marked set
$\mathcal M$,
\[
 \sum_{D\in\mathcal F}\prod_{v\in D}w_v
 \le \exp\left\{\sum_{v\in\mathcal M}w_v+
 \sum_vw_v\left(\prod_{u\sim v}(1+w_u)-1\right)\right\}.
\]
\end{lemma}
\noindent The complete proof is given in Section~\ref{p14:comp:certificate-generator}.
